\chapter*{Справочник обозначений}\markright{Обозначения}
\addcontentsline{toc}{chapter}{Справочник обозначений}
\begingroup\small\renewcommand{\arraystretch}{1.35}\setlength{\tabcolsep}{4pt}\arrayrulecolor{Teal}\setlength{\arrayrulewidth}{.35pt}
\Needspace{7\baselineskip}\section*{Начертания символов}
\begin{longtable}{|>{\raggedright\arraybackslash}p{.21\textwidth}|>{\raggedright\arraybackslash}p{.27\textwidth}|>{\raggedright\arraybackslash}p{.46\textwidth}|}
\hline\textbf{Обозначение}&\textbf{Как читать}&\textbf{Значение}\\\hline\endfirsthead
\hline\textbf{Обозначение}&\textbf{Как читать}&\textbf{Значение}\\\hline\endhead
\hline\endfoot
\hline\endlastfoot
\(\displaystyle x\)&математический курсив&Обычное начертание переменных; смысл буквы задаётся в тексте.\\\hline
\(\displaystyle \mathrm{x}\)&прямое начертание&Так оформляют названия операторов и единицы измерения; начертание само по себе не задаёт значение.\\\hline
\(\displaystyle \symbf v=\vec v\)&полужирное начертание; стрелка&Две распространённые записи вектора.\\\hline
\(\displaystyle \mathbb R\)&двойное начертание&Распространено для числовых множеств; здесь показан знак действительных чисел.\\\hline
\(\displaystyle \mathcal A\)&каллиграфическое начертание&Часто обозначает семейство множеств или другой специально определённый объект.\\\hline
\(\displaystyle \mathfrak g\)&фрактура&Используется в алгебре для отличия одних объектов от других; смысл устанавливает определение.\\\hline
\end{longtable}
\Needspace{7\baselineskip}\section*{Числа и числовые множества}
\begin{longtable}{|>{\raggedright\arraybackslash}p{.21\textwidth}|>{\raggedright\arraybackslash}p{.27\textwidth}|>{\raggedright\arraybackslash}p{.46\textwidth}|}
\hline\textbf{Обозначение}&\textbf{Как читать}&\textbf{Значение}\\\hline\endfirsthead
\hline\textbf{Обозначение}&\textbf{Как читать}&\textbf{Значение}\\\hline\endhead
\hline\endfoot
\hline\endlastfoot
\(\displaystyle 0\)&ноль&Нейтральный элемент сложения: \(a+0=a\).\\\hline
\(\displaystyle 1\)&единица&Нейтральный элемент умножения: \(a\cdot1=a\).\\\hline
\(\displaystyle \mathbb N\)&натуральные числа&\(\{1,2,3,\ldots\}\).\\\hline
\(\displaystyle \mathbb N_0\)&натуральные числа с нулём&\(\{0,1,2,3,\ldots\}\).\\\hline
\(\displaystyle \mathbb Z\)&целые числа&\(\{\ldots,-2,-1,0,1,2,\ldots\}\).\\\hline
\(\displaystyle \mathbb Q\)&рациональные числа&Числа \(\frac pq\), где \(p\in\mathbb Z\), \(q\in\mathbb Z\), \(q\ne0\).\\\hline
\(\displaystyle \mathbb R\)&действительные числа&Числа числовой прямой, включая рациональные и иррациональные.\\\hline
\(\displaystyle \mathbb C\)&комплексные числа&Числа вида \(a+bi\), где \(a,b\in\mathbb R\).\\\hline
\(\displaystyle \mathbb H\)&кватернионы&Вещественная алгебра размерности 4 с некоммутативным ассоциативным умножением.\\\hline
\(\displaystyle \mathbb O\)&октонионы&Вещественная алгебра размерности 8; умножение не ассоциативно.\\\hline
\(\displaystyle \mathbb K\)&поле скаляров&В главах о векторных пространствах: \(\mathbb R\) или \(\mathbb C\), по условию.\\\hline
\(\displaystyle \mathbb R^n\)&эр в степени эн&Множество упорядоченных наборов из \(n\) действительных координат.\\\hline
\end{longtable}
\Needspace{7\baselineskip}\section*{Действия и сравнения чисел}
\begin{longtable}{|>{\raggedright\arraybackslash}p{.21\textwidth}|>{\raggedright\arraybackslash}p{.27\textwidth}|>{\raggedright\arraybackslash}p{.46\textwidth}|}
\hline\textbf{Обозначение}&\textbf{Как читать}&\textbf{Значение}\\\hline\endfirsthead
\hline\textbf{Обозначение}&\textbf{Как читать}&\textbf{Значение}\\\hline\endhead
\hline\endfoot
\hline\endlastfoot
\(\displaystyle +\)&плюс&Сложение: \(2+3=5\).\\\hline
\(\displaystyle -\)&минус&Вычитание или противоположное число: \(5-2=3\), \(-a\).\\\hline
\(\displaystyle \cdot\)&умножить&Умножение: \(ab=a\cdot b\).\\\hline
\(\displaystyle :\)&разделить&Деление: \(6:3=2\); делитель ненулевой.\\\hline
\(\displaystyle \frac ab\)&дробь а на б&Частное; знаменатель \(b\ne0\).\\\hline
\(\displaystyle =\)&равно&Равенство двух выражений.\\\hline
\(\displaystyle \ne\)&не равно&Выражения имеют разные значения.\\\hline
\(\displaystyle <\)&меньше&Строгое сравнение.\\\hline
\(\displaystyle >\)&больше&Строгое сравнение.\\\hline
\(\displaystyle \le\)&меньше или равно&Нестрогое сравнение; равенство допускается.\\\hline
\(\displaystyle \ge\)&больше или равно&Нестрогое сравнение; равенство допускается.\\\hline
\(\displaystyle \approx\)&приблизительно равно&Например, \(\sqrt2\approx1{,}414\).\\\hline
\(\displaystyle \pm\)&плюс или минус&\(x=\pm2\) означает \(x=-2\) или \(x=2\).\\\hline
\(\displaystyle a^n\)&а в степени эн&Степень с натуральным показателем; \(a^0=1\) при \(a\ne0\).\\\hline
\(\displaystyle a^{-n}\)&а в степени минус эн&При \(n\in\mathbb N\), \(a\ne0\): \(a^{-n}=\frac1{a^n}\).\\\hline
\(\displaystyle \sqrt a\)&квадратный корень из а&Неотрицательное число с квадратом \(a\); требуется \(a\ge0\).\\\hline
\(\displaystyle \sqrt[n]a\)&корень энной степени&Для чётного \(n\) берётся неотрицательный корень и требуется \(a\ge0\).\\\hline
\(\displaystyle |a|\)&модуль а&Расстояние от \(a\) до нуля.\\\hline
\end{longtable}
\Needspace{7\baselineskip}\section*{Делимость и сравнения по модулю}
\begin{longtable}{|>{\raggedright\arraybackslash}p{.21\textwidth}|>{\raggedright\arraybackslash}p{.27\textwidth}|>{\raggedright\arraybackslash}p{.46\textwidth}|}
\hline\textbf{Обозначение}&\textbf{Как читать}&\textbf{Значение}\\\hline\endfirsthead
\hline\textbf{Обозначение}&\textbf{Как читать}&\textbf{Значение}\\\hline\endhead
\hline\endfoot
\hline\endlastfoot
\(\displaystyle a\mathrel{\vdots}b\)&а делится на б нацело&\(a=bk\) для целого \(k\); \(b\ne0\).\\\hline
\(\displaystyle a\mid b\)&а делит б&\(b=ak\) для целого \(k\); порядок обратен записи с тремя точками.\\\hline
\(\displaystyle a\nmid b\)&а не делит б&Число \(b\) не кратно \(a\).\\\hline
\(\displaystyle \gcd(a,b)\)&наибольший общий делитель&НОД; английское greatest common divisor.\\\hline
\(\displaystyle \operatorname{lcm}(a,b)\)&наименьшее общее кратное&НОК; английское least common multiple.\\\hline
\(\displaystyle a\equiv b\pmod m\)&а сравнимо с б по модулю эм&\(m\mid(a-b)\), где \(m\) — положительное целое.\\\hline
\(\displaystyle [a]_m\)&класс а по модулю эм&Все целые числа \(a+mk\), \(k\in\mathbb Z\).\\\hline
\(\displaystyle \mathbb Z/m\mathbb Z\)&целые по модулю эм&Множество классов вычетов по модулю \(m\).\\\hline
\end{longtable}
\Needspace{7\baselineskip}\section*{Множества и пары}
\begin{longtable}{|>{\raggedright\arraybackslash}p{.21\textwidth}|>{\raggedright\arraybackslash}p{.27\textwidth}|>{\raggedright\arraybackslash}p{.46\textwidth}|}
\hline\textbf{Обозначение}&\textbf{Как читать}&\textbf{Значение}\\\hline\endfirsthead
\hline\textbf{Обозначение}&\textbf{Как читать}&\textbf{Значение}\\\hline\endhead
\hline\endfoot
\hline\endlastfoot
\(\displaystyle x\in A\)&икс принадлежит а&\(x\) — элемент множества \(A\).\\\hline
\(\displaystyle x\notin A\)&икс не принадлежит а&\(x\) не является элементом \(A\).\\\hline
\(\displaystyle \{\ldots\}\)&множество перечисленных элементов&Фигурные скобки: \(\{1,2,3\}\).\\\hline
\(\displaystyle \varnothing\)&пустое множество&Множество без элементов.\\\hline
\(\displaystyle A\subseteq B\)&а является подмножеством б&Каждый элемент \(A\) принадлежит \(B\); равенство допускается.\\\hline
\(\displaystyle A\subset B\)&а — собственное подмножество б&\(A\subseteq B\), но \(A\ne B\).\\\hline
\(\displaystyle A\cup B\)&объединение а и б&Элементы хотя бы одного из множеств.\\\hline
\(\displaystyle A\cap B\)&пересечение а и б&Элементы, принадлежащие обоим множествам.\\\hline
\(\displaystyle A\setminus B\)&разность а и б&Элементы \(A\), не принадлежащие \(B\).\\\hline
\(\displaystyle A^c\)&дополнение а&\(U\setminus A\); зависит от выбранного универсума.\\\hline
\(\displaystyle A\triangle B\)&симметрическая разность&Элементы ровно одного из двух множеств.\\\hline
\(\displaystyle A\times B\)&декартово произведение&Все упорядоченные пары с первой координатой из \(A\), второй из \(B\).\\\hline
\(\displaystyle (a,b)\)&упорядоченная пара&Порядок координат важен.\\\hline
\(\displaystyle \mathcal P(A)=2^A\)&множество всех подмножеств а&Все множества \(B\), для которых \(B\subseteq A\).\\\hline
\(\displaystyle |A|\)&число элементов а&Для конечного множества; смысл отличается от модуля числа.\\\hline
\end{longtable}
\Needspace{7\baselineskip}\section*{Промежутки числовой прямой}
\begin{longtable}{|>{\raggedright\arraybackslash}p{.21\textwidth}|>{\raggedright\arraybackslash}p{.27\textwidth}|>{\raggedright\arraybackslash}p{.46\textwidth}|}
\hline\textbf{Обозначение}&\textbf{Как читать}&\textbf{Значение}\\\hline\endfirsthead
\hline\textbf{Обозначение}&\textbf{Как читать}&\textbf{Значение}\\\hline\endhead
\hline\endfoot
\hline\endlastfoot
\(\displaystyle (a;b)\)&открытый интервал&\(a<x<b\); концы исключены.\\\hline
\(\displaystyle [a;b]\)&отрезок&\(a\le x\le b\); оба конца включены.\\\hline
\(\displaystyle [a;b)\)&полуинтервал&\(a\le x<b\); левый конец включён, правый исключён.\\\hline
\(\displaystyle (a;b]\)&полуинтервал&\(a<x\le b\); левый конец исключён, правый включён.\\\hline
\(\displaystyle \infty\)&бесконечность&Не действительное число; у бесконечной границы промежутка ставят круглую скобку.\\\hline
\end{longtable}
\Needspace{7\baselineskip}\section*{Логика и доказательства}
\begin{longtable}{|>{\raggedright\arraybackslash}p{.21\textwidth}|>{\raggedright\arraybackslash}p{.27\textwidth}|>{\raggedright\arraybackslash}p{.46\textwidth}|}
\hline\textbf{Обозначение}&\textbf{Как читать}&\textbf{Значение}\\\hline\endfirsthead
\hline\textbf{Обозначение}&\textbf{Как читать}&\textbf{Значение}\\\hline\endhead
\hline\endfoot
\hline\endlastfoot
\(\displaystyle \forall\)&для любого&Квантор всеобщности.\\\hline
\(\displaystyle \exists\)&существует&Квантор существования.\\\hline
\(\displaystyle \to,\ \Rightarrow\)&импликация&Если \(P\) истинно, то \(Q\) истинно.\\\hline
\(\displaystyle \leftrightarrow,\ \Longleftrightarrow\)&тогда и только тогда&Каждое утверждение следует из другого.\\\hline
\(\displaystyle \neg P\)&не пэ&Отрицание утверждения \(P\).\\\hline
\(\displaystyle P\land Q\)&пэ и ку&Оба утверждения истинны.\\\hline
\(\displaystyle P\lor Q\)&пэ или ку&Хотя бы одно утверждение истинно.\\\hline
\(\displaystyle \blacksquare\)&конец доказательства&В конце доказательства; в отдельных заданиях пустой квадрат может обозначать неизвестную цифру.\\\hline
\(\displaystyle \ast\)&звёздочка&Задача повышенной сложности.\\\hline
\(\displaystyle P\oplus Q\)&исключающее или&Истинно, когда истинно ровно одно из двух утверждений.\\\hline
\(\displaystyle P\uparrow Q\)&штрих Шеффера&Отрицание конъюнкции: \(\neg(P\land Q)\).\\\hline
\(\displaystyle P\downarrow Q\)&стрелка Пирса&Отрицание дизъюнкции: \(\neg(P\lor Q)\).\\\hline
\(\displaystyle \top\)&истина&Логическая константа истины.\\\hline
\(\displaystyle \bot\)&ложь&Логическая константа лжи.\\\hline
\(\displaystyle \Gamma\vdash A\)&из гаммы выводится а&Синтаксическая выводимость в указанной системе правил.\\\hline
\(\displaystyle \Gamma\models A\)&из гаммы семантически следует а&При каждой оценке, делающей все посылки истинными, заключение истинно.\\\hline
\(\displaystyle A\equiv B\)&формулы а и б равносильны&Формулы имеют одно истинностное значение при каждой оценке.\\\hline
\(\displaystyle v\models A\)&оценка вэ удовлетворяет а&Формула \(A\) истинна при оценке \(v\).\\\hline
\end{longtable}
\Needspace{7\baselineskip}\section*{Суммы, комбинаторика и последовательности}
\begin{longtable}{|>{\raggedright\arraybackslash}p{.21\textwidth}|>{\raggedright\arraybackslash}p{.27\textwidth}|>{\raggedright\arraybackslash}p{.46\textwidth}|}
\hline\textbf{Обозначение}&\textbf{Как читать}&\textbf{Значение}\\\hline\endfirsthead
\hline\textbf{Обозначение}&\textbf{Как читать}&\textbf{Значение}\\\hline\endhead
\hline\endfoot
\hline\endlastfoot
\(\displaystyle S(n)\)&следующее за эн&Операция следования в аксиомах Пеано; \(S(0)=1\).\\\hline
\(\displaystyle \sum_{k=1}^{n}a_k\)&сумма по ка от одного до эн&\(a_1+\cdots+a_n\).\\\hline
\(\displaystyle \prod_{k=1}^{n}a_k\)&произведение по ка&\(a_1\cdots a_n\).\\\hline
\(\displaystyle n!\)&эн факториал&\(1\cdot2\cdots n\) при \(n\ge1\); отдельно \(0!=1\).\\\hline
\(\displaystyle P_n\)&число перестановок&Перестановки \(n\) различных элементов: \(n!\).\\\hline
\(\displaystyle A_n^k\)&число размещений&Выбор \(k\) из \(n\) различных объектов с учётом порядка.\\\hline
\(\displaystyle \binom nk=C_n^k\)&число сочетаний из эн по ка&Выбор \(k\) из \(n\) без учёта порядка.\\\hline
\(\displaystyle a_n\)&энный член последовательности&Значение последовательности с номером \(n\).\\\hline
\(\displaystyle S_n\)&энная частичная сумма&Сумма первых \(n\) членов последовательности.\\\hline
\(\displaystyle \lim_{n\to\infty}a_n\)&предел последовательности&Число, к которому члены становятся сколь угодно близки при больших \(n\).\\\hline
\(\displaystyle \sup A\)&точная верхняя грань а&Наименьшая из верхних границ множества \(A\), если она существует.\\\hline
\(\displaystyle \inf A\)&точная нижняя грань а&Наибольшая из нижних границ непустого множества \(A\), если она существует.\\\hline
\end{longtable}
\Needspace{7\baselineskip}\section*{Функции, экспонента и логарифмы}
\begin{longtable}{|>{\raggedright\arraybackslash}p{.21\textwidth}|>{\raggedright\arraybackslash}p{.27\textwidth}|>{\raggedright\arraybackslash}p{.46\textwidth}|}
\hline\textbf{Обозначение}&\textbf{Как читать}&\textbf{Значение}\\\hline\endfirsthead
\hline\textbf{Обозначение}&\textbf{Как читать}&\textbf{Значение}\\\hline\endhead
\hline\endfoot
\hline\endlastfoot
\(\displaystyle f:X\to Y\)&эф из икс в игрек&Отображение с областью определения \(X\) и множеством назначения \(Y\).\\\hline
\(\displaystyle x\mapsto f(x)\)&икс переходит в эф от икс&Правило, сопоставляющее аргументу значение функции.\\\hline
\(\displaystyle f(x)\)&эф от икс&Значение функции при аргументе \(x\).\\\hline
\(\displaystyle f\circ g\)&композиция эф и жэ&\((f\circ g)(x)=f(g(x))\); сначала выполняется \(g\).\\\hline
\(\displaystyle f^{-1}\)&обратная функция&Обратное отображение; обычно это не \(\frac1f\).\\\hline
\(\displaystyle e\)&число е&Основание натурального логарифма; \(e=\lim_{n\to\infty}(1+\frac1n)^n\).\\\hline
\(\displaystyle e^x=\exp x\)&экспонента от икс&Показательная функция с основанием \(e\).\\\hline
\(\displaystyle \ln x=\log_e x\)&натуральный логарифм&Логарифм по основанию \(e\); \(x>0\).\\\hline
\(\displaystyle \log_a x\)&логарифм икс по основанию а&\(a>0\), \(a\ne1\), \(x>0\).\\\hline
\end{longtable}
\Needspace{7\baselineskip}\section*{Тригонометрические и гиперболические функции}
\begin{longtable}{|>{\raggedright\arraybackslash}p{.21\textwidth}|>{\raggedright\arraybackslash}p{.27\textwidth}|>{\raggedright\arraybackslash}p{.46\textwidth}|}
\hline\textbf{Обозначение}&\textbf{Как читать}&\textbf{Значение}\\\hline\endfirsthead
\hline\textbf{Обозначение}&\textbf{Как читать}&\textbf{Значение}\\\hline\endhead
\hline\endfoot
\hline\endlastfoot
\(\displaystyle \pi\)&число пи&Отношение длины окружности к её диаметру.\\\hline
\(\displaystyle \sin x\)&синус икс&Вертикальная координата на единичной окружности.\\\hline
\(\displaystyle \cos x\)&косинус икс&Горизонтальная координата на единичной окружности.\\\hline
\(\displaystyle \operatorname{tg} x\)&тангенс икс&\(\frac{\sin x}{\cos x}\) при \(\cos x\ne0\).\\\hline
\(\displaystyle \arcsin x\)&арксинус икс&Обратная к синусу на \([-\frac\pi2;\frac\pi2]\).\\\hline
\(\displaystyle \arccos x\)&арккосинус икс&Обратная к косинусу на \([0;\pi]\).\\\hline
\(\displaystyle \operatorname{arctg} x\)&арктангенс икс&Обратная к тангенсу на \((-\frac\pi2;\frac\pi2)\).\\\hline
\(\displaystyle \operatorname{sh} x\)&гиперболический синус&\(\frac{e^x-e^{-x}}2\).\\\hline
\(\displaystyle \operatorname{ch} x\)&гиперболический косинус&\(\frac{e^x+e^{-x}}2\).\\\hline
\(\displaystyle \operatorname{ctg}x\)&котангенс икс&\(\frac{\cos x}{\sin x}\); \(\sin x\ne0\).\\\hline
\(\displaystyle \operatorname{th}x\)&гиперболический тангенс&\(\frac{\operatorname{sh}x}{\operatorname{ch}x}\); определён при всех действительных \(x\).\\\hline
\(\displaystyle \operatorname{cth}x\)&гиперболический котангенс&\(\frac{\operatorname{ch}x}{\operatorname{sh}x}\); \(x\ne0\).\\\hline
\end{longtable}
\Needspace{7\baselineskip}\section*{Векторы и матрицы}
\begin{longtable}{|>{\raggedright\arraybackslash}p{.21\textwidth}|>{\raggedright\arraybackslash}p{.27\textwidth}|>{\raggedright\arraybackslash}p{.46\textwidth}|}
\hline\textbf{Обозначение}&\textbf{Как читать}&\textbf{Значение}\\\hline\endfirsthead
\hline\textbf{Обозначение}&\textbf{Как читать}&\textbf{Значение}\\\hline\endhead
\hline\endfoot
\hline\endlastfoot
\(\displaystyle \lVert v\rVert\)&длина вектора вэ&Евклидова длина, например \(\sqrt{x^2+y^2}\) для \(v=(x,y)\).\\\hline
\(\displaystyle a_{ij}\)&элемент а с индексами и, жэ&Элемент матрицы в строке \(i\), столбце \(j\).\\\hline
\(\displaystyle AB\)&произведение матриц а и б&Определено, если число столбцов \(A\) равно числу строк \(B\).\\\hline
\(\displaystyle I_n\)&единичная матрица порядка эн&На главной диагонали стоят единицы, остальные элементы равны нулю.\\\hline
\(\displaystyle \det A\)&определитель матрицы а&Для \(A=\begin{pmatrix}a&b\\c&d\end{pmatrix}\) равен \(ad-bc\).\\\hline
\(\displaystyle A^{-1}\)&обратная матрица&\(AA^{-1}=A^{-1}A=I\); существует не у всякой квадратной матрицы.\\\hline
\(\displaystyle \symbf 0\)&нулевой вектор&Вектор, все координаты которого равны нулю.\\\hline
\(\displaystyle \langle u,v\rangle=u\cdot v\)&скалярное произведение векторов&Евклидово скалярное произведение в \(\mathbb R^n\): \(\sum_{k=1}^nu_kv_k\).\\\hline
\(\displaystyle \operatorname{rank} A\)&ранг матрицы а&Число ведущих элементов в ступенчатом виде матрицы.\\\hline
\(\displaystyle \operatorname{span}S\)&линейная оболочка&Все конечные линейные комбинации векторов из указанного набора.\\\hline
\(\displaystyle \dim V\)&размерность вэ&Число векторов в базисе конечномерного пространства.\\\hline
\(\displaystyle [v]_B\)&координаты вэ в базисе бэ&Столбец коэффициентов разложения вектора по указанному базису.\\\hline
\(\displaystyle \ker T\)&ядро тэ&Векторы \(v\), для которых \(T(v)=0\).\\\hline
\(\displaystyle \operatorname{im}T\)&образ тэ&Все значения \(T(v)\).\\\hline
\(\displaystyle \mathcal P_m\)&пространство многочленов&Многочлены степени не выше \(m\), включая нулевой многочлен.\\\hline
\(\displaystyle A^{\mathrm T}\)&транспонированная матрица&Строки и столбцы меняются местами.\\\hline
\(\displaystyle \operatorname{diag}(\ldots)\)&диагональная матрица&Указанные элементы на диагонали, нули вне диагонали.\\\hline
\(\displaystyle \chi_A(t)\)&характеристический многочлен&В главе 27 используется соглашение \(\det(tI-A)\); корни — собственные значения.\\\hline
\(\displaystyle E_\lambda\)&собственное подпространство&\(\ker(A-\lambda I)\); включает нулевой вектор.\\\hline
\end{longtable}
\Needspace{7\baselineskip}\section*{Предел функции и производная}
\begin{longtable}{|>{\raggedright\arraybackslash}p{.21\textwidth}|>{\raggedright\arraybackslash}p{.27\textwidth}|>{\raggedright\arraybackslash}p{.46\textwidth}|}
\hline\textbf{Обозначение}&\textbf{Как читать}&\textbf{Значение}\\\hline\endfirsthead
\hline\textbf{Обозначение}&\textbf{Как читать}&\textbf{Значение}\\\hline\endhead
\hline\endfoot
\hline\endlastfoot
\(\displaystyle x\to a\)&икс стремится к а&Аргумент приближается к точке \(a\).\\\hline
\(\displaystyle x\to a-\)&икс стремится к а слева&Рассматриваются значения \(x<a\).\\\hline
\(\displaystyle x\to a+\)&икс стремится к а справа&Рассматриваются значения \(x>a\).\\\hline
\(\displaystyle \lim_{x\to a}f(x)\)&предел эф от икс при икс к а&Предел в точке; значение \(f(a)\) может отсутствовать или отличаться от предела.\\\hline
\(\displaystyle f'(a)\)&производная эф в точке а&Предел \(\lim_{h\to0}\frac{f(a+h)-f(a)}h\), если он существует и конечен.\\\hline
\(\displaystyle y-f(a)=f'(a)(x-a)\)&уравнение касательной&Касательная к графику дифференцируемой функции в точке с абсциссой \(a\).\\\hline
\end{longtable}
\Needspace{7\baselineskip}\section*{Алгебраические структуры}
\begin{longtable}{|>{\raggedright\arraybackslash}p{.21\textwidth}|>{\raggedright\arraybackslash}p{.27\textwidth}|>{\raggedright\arraybackslash}p{.46\textwidth}|}
\hline\textbf{Обозначение}&\textbf{Как читать}&\textbf{Значение}\\\hline\endfirsthead
\hline\textbf{Обозначение}&\textbf{Как читать}&\textbf{Значение}\\\hline\endhead
\hline\endfoot
\hline\endlastfoot
\(\displaystyle a*b\)&а звёздочка б&Результат заданной бинарной операции; это не обязательно обычное умножение.\\\hline
\(\displaystyle e_G\)&нейтральный элемент группы жэ&Элемент, не меняющий другие элементы при операции.\\\hline
\(\displaystyle a^{-1}\)&обратный к а&В группе: \(a*a^{-1}=a^{-1}*a=e\).\\\hline
\(\displaystyle \operatorname{id}_X\)&тождественное отображение икс&Оставляет каждый элемент \(X\) на месте.\\\hline
\(\displaystyle \operatorname{Sym}(X)\)&группа перестановок икс&Все биекции множества \(X\) в себя с операцией композиции.\\\hline
\(\displaystyle S_n\)&симметрическая группа степени эн&Все перестановки \(\{1,\ldots,n\}\); число элементов этой группы равно \(n!\).\\\hline
\(\displaystyle \mathrm M_2(\mathbb R)\)&матрицы два на два над эр&Кольцо действительных матриц размера \(2\times2\).\\\hline
\(\displaystyle \mathbb F_p\)&поле из пэ элементов&При простом \(p\): поле вычетов \(\mathbb Z/p\mathbb Z\).\\\hline
\end{longtable}
\Needspace{7\baselineskip}\section*{Комплексные числа}
\begin{longtable}{|>{\raggedright\arraybackslash}p{.21\textwidth}|>{\raggedright\arraybackslash}p{.27\textwidth}|>{\raggedright\arraybackslash}p{.46\textwidth}|}
\hline\textbf{Обозначение}&\textbf{Как читать}&\textbf{Значение}\\\hline\endfirsthead
\hline\textbf{Обозначение}&\textbf{Как читать}&\textbf{Значение}\\\hline\endhead
\hline\endfoot
\hline\endlastfoot
\(\displaystyle i\)&мнимая единица&Комплексное число, для которого \(i^2=-1\).\\\hline
\(\displaystyle a+bi\)&а плюс бэ и&Алгебраическая форма комплексного числа.\\\hline
\(\displaystyle \operatorname{Re}z\)&действительная часть зет&Число \(a\) в записи \(z=a+bi\).\\\hline
\(\displaystyle \operatorname{Im}z\)&мнимая часть зет&Действительное число \(b\) в записи \(z=a+bi\); множитель \(i\) в мнимую часть не входит.\\\hline
\(\displaystyle \overline z\)&зет с чертой&Комплексное сопряжение: при \(z=a+bi\) получаем \(\overline z=a-bi\).\\\hline
\(\displaystyle |z|\)&модуль комплексного числа&При \(z=a+bi\): \(\sqrt{a^2+b^2}\). Расстояние до начала координат.\\\hline
\(\displaystyle \arg z\)&аргументы зет&Все углы направления ненулевого \(z\); различаются на \(2\pi k\), \(k\in\mathbb Z\). У нуля аргумент не определён.\\\hline
\(\displaystyle \operatorname{Arg}z\)&главный аргумент зет&Для \(z\ne0\): единственный аргумент из \((-\pi;\pi]\); это принятое в книге соглашение.\\\hline
\(\displaystyle r(\cos\varphi+i\sin\varphi)\)&тригонометрическая форма&Для \(z\ne0\) здесь \(r=|z|>0\), а \(\varphi\) — один из аргументов.\\\hline
\(\displaystyle e^{i\varphi}\)&е в степени и фи&Экспоненциальная запись \(\cos\varphi+i\sin\varphi\) для действительного \(\varphi\).\\\hline
\end{longtable}
\Needspace{7\baselineskip}\section*{Кватернионы и октонионы}
\begin{longtable}{|>{\raggedright\arraybackslash}p{.21\textwidth}|>{\raggedright\arraybackslash}p{.27\textwidth}|>{\raggedright\arraybackslash}p{.46\textwidth}|}
\hline\textbf{Обозначение}&\textbf{Как читать}&\textbf{Значение}\\\hline\endfirsthead
\hline\textbf{Обозначение}&\textbf{Как читать}&\textbf{Значение}\\\hline\endhead
\hline\endfoot
\hline\endlastfoot
\(\displaystyle j\)&кватернионная единица жэ&\(j^2=-1\), \(ij=k\), \(ji=-k\).\\\hline
\(\displaystyle k\)&кватернионная единица ка&\(k=ij\), \(k^2=-1\).\\\hline
\(\displaystyle N(q)=|q|^2\)&квадрат нормы&Сумма квадратов действительных координат кватерниона или октониона.\\\hline
\(\displaystyle e_0,\ldots,e_7\)&базис октонионов&Восемь базисных элементов в выбранном соглашении; \(e_0=1\).\\\hline
\(\displaystyle [x,y,z]\)&ассоциатор&Разность \((xy)z-x(yz)\); измеряет зависимость от расстановки скобок.\\\hline
\end{longtable}
\Needspace{7\baselineskip}\section*{Метрики и окрестности}
\begin{longtable}{|>{\raggedright\arraybackslash}p{.21\textwidth}|>{\raggedright\arraybackslash}p{.27\textwidth}|>{\raggedright\arraybackslash}p{.46\textwidth}|}
\hline\textbf{Обозначение}&\textbf{Как читать}&\textbf{Значение}\\\hline\endfirsthead
\hline\textbf{Обозначение}&\textbf{Как читать}&\textbf{Значение}\\\hline\endhead
\hline\endfoot
\hline\endlastfoot
\(\displaystyle d(x,y)\)&расстояние между икс и игрек&Значение выбранной метрики; её правило задаётся в условии.\\\hline
\(\displaystyle d_1(x,y)\)&метрика суммы модулей&\(\sum_j|x_j-y_j|\).\\\hline
\(\displaystyle d_2(x,y)\)&евклидова метрика&\(\sqrt{\sum_j(x_j-y_j)^2}\).\\\hline
\(\displaystyle d_\infty(x,y)\)&метрика максимума&\(\max_j|x_j-y_j|\).\\\hline
\(\displaystyle B_r(a)\)&открытый шар&Точки \(x\), для которых \(d(x,a)<r\), где \(r>0\).\\\hline
\(\displaystyle (x_{n_k})\)&подпоследо\-вательность&Члены с номерами \(n_1<n_2<\cdots\).\\\hline
\end{longtable}
\Needspace{7\baselineskip}\section*{Интегралы и ряды}
\begin{longtable}{|>{\raggedright\arraybackslash}p{.21\textwidth}|>{\raggedright\arraybackslash}p{.27\textwidth}|>{\raggedright\arraybackslash}p{.46\textwidth}|}
\hline\textbf{Обозначение}&\textbf{Как читать}&\textbf{Значение}\\\hline\endfirsthead
\hline\textbf{Обозначение}&\textbf{Как читать}&\textbf{Значение}\\\hline\endhead
\hline\endfoot
\hline\endlastfoot
\(\displaystyle \int_a^b f(x)\,dx\)&определённый интеграл&Интеграл от нижней границы до верхней.\\\hline
\(\displaystyle \int f(x)\,dx\)&неопределённый интеграл&Семейство первообразных на указанном интервале.\\\hline
\(\displaystyle [F(x)]_a^b\)&подстановка пределов&Разность \(F(b)-F(a)\).\\\hline
\(\displaystyle \|P\|\)&мелкость разбиения&Наибольшая длина его частей; смысл отличается от нормы вектора.\\\hline
\(\displaystyle L(f,P)\)&нижняя сумма Дарбу&\(\sum_i m_i(x_i-x_{i-1})\); используются инфимумы на частях.\\\hline
\(\displaystyle U(f,P)\)&верхняя сумма Дарбу&\(\sum_i M_i(x_i-x_{i-1})\); используются супремумы на частях.\\\hline
\(\displaystyle \sum_{n=1}^\infty a_n\)&числовой ряд&Предел конечных частичных сумм, если он существует и конечен.\\\hline
\(\displaystyle \textstyle\sum_{n=0}^\infty c_n(x-a)^n\)&степенной ряд&Для каждого фиксированного аргумента получается числовой ряд.\\\hline
\end{longtable}
\endgroup
